%%%%

%%%   Самарский государственный технический университет

%%%   Кафедра прикладной математики и информатики

%%%

%%%   Преамбула для статьи в журнал "Вестник Самарского государственного

%%%   технического университета. Серия: «Физико-математические науки»"

%%%   Ориентирована под MikTEx 2.4 for Win32

%%%

%%%   Хотя сам журнал собирается под GNU/Linux на дистрибутиве TeXLive



\documentclass[11pt,a4paper,twoside]{article}



\usepackage[cp1251]{inputenc}

\usepackage[T2A]{fontenc}

\usepackage[english,russian]{babel}

\usepackage{samgtubib}

\usepackage[dvips]{graphicx}

\usepackage[multiple]{footmisc}



\usepackage{samgtu}

\renewcommand{\VestnikYear}{2009} % Год издания Вестника

\renewcommand{\VestnikNumber}{2\,(19)} % Номер Вестника

\renewcommand{\VestnikMonth}{Ноябрь} % Месяц выхода Вестника

\renewcommand{\VestnikData}{01.11.09} % Дата подписания Вестника в печать

\renewcommand{\VestnikUPL}{26,64} % Усл. печ. л.

\renewcommand{\VestnikUIL}{25,73} % Уч.-изд. л.

\renewcommand{\VestnikRegNumber}{479/09} % Регистрационный номер

\renewcommand{\VestnikOrder}{994} % Номер заказа







%

%

%\begin{document}



\renewcommand{\firstpage}{8}

\renewcommand{\lastpage}{25}









%%%%%%%%%%%%%%%%%%%%%%%%%%%%%%%%%%%%%%%%%%%%%%%%%%%%%%%%%%%%%%%%%%%%%%%%%%%%%%%%%%%%



\udk{517.956.6}

\msc{35M10}



\title{О некоторых задачах сопряжения уравнений параболического и~гиперболического типов с~интегро-диф\-фе\-рен\-циаль\-ными условиями

на границе раздела областей}{О некоторых задачах сопряжения уравнений параболического и~гиперболического типов \dots}



\author{В.\,А.~Елеев, А.\,Х.~Балкизова}

{Е\,л\,е\,е\,в ~В.\,А.,\ Б\,а\,л\,к\,и\,з\,о\,в\,а ~А.\,Х.}



\institute{Кабардино-Балкарский государственный университет им.~Х.~М.~Бербекова,\\

360004, Нальчик, ул.~Чернышевского, 173.}



\email{E-mails}{eleev@yandex.ru, alena-balkizova@rambler.ru}



\otherinf{\fullauthor{Валерий Абдурахманович Елеев}\regalia{д.ф.-м.н., профессор}\position{зав. кафедрой}

\dept{каф. теории функций и функционального анализа} 

\fullauthor{Алёна Хамидбиевна Балкизова} \position{аспирант} \dept{каф. теории функций и функционального анализа}  }



\keywords{интегральное уравнение Вольтерра, задачи сопряжения, функции Бесселя}



\workabstract{Доказана однозначная разрешимость задач сопряжения гиперболического и~параболического уравнений

в~конечных областях методом эквивалентной редукции к~интегральному уравнению Вольтерра второго рода.}



\engtitle{On some conjugation problems of parabolic and hyperbolic equations with integro-differential conditions on the separating boundary}



\engauthor{V.\,A.~Eleev, A.\,H.~Balkizova}{E\,l\,e\,e\,v~V.\,A.,

B\,a\,l\,k\,i\,z\,o\,v\,a~A.\,H.}



\enginstitute{Kabardino-Balkar State University,\\

173, Chernyshevskogo st., Nal'chik, 360004.}



\otherenginf{\fullauthor{Valeriy A. Eleev} \regalia{Dr.\,Sci. (Phys. \& Math.)} \position{Head of Dept.}

\dept{Dept. of Function Theory and Functional Analysis}

\fullauthor{Alyona H. Balkizova}

\position{Postgraduate Student} \dept{Dept. of Function Theory and Functional Analysis}}



\engabstract{The one-valued solvability of the conjugation problems of parabolic and hyperbolic equations in finite domains was proved by the method of equivalent reduction to Volterra integral equation of the second kind.

}



\engkeywords{Volterra integral equation, conjugation problems, Bessel functions}



\date{18/XI/2010}

\revisiondate{25/VIII/2011}



%%%%%%%%%%%%%%%%%%%%%%%%%%%%%%%%%%%%%%%%%%%%%%%%%%%%%%%%%%%%%%%%%%%%%%%%%%%%%%%%%%%%



\maketitle \vspace{-3mm}



Необходимость рассмотрения сопряжения, когда на одной части области задано уравнение параболического типа, а~на другой "--- уравнение гиперболического типа, была впервые высказана И.\,М.~Гельфандом в 1959~г. [\citen{eleevbalk:bib:Gel}].

К~задаче сопряжения приводит изучение электрических колебаний в~проводах.

Такого рода задачи встречаются также при изучении движения жидкости в~канале, окруженной пористой средой, в~теории

распространения электромагнитных полей и~в~ряде других областей физики.

Так, в~канале гидродинамическое давление жидкости удовлетворяет волновому уравнению, а~в~пористой среде "--- уравнению фильтрации,

которое в~данном случае совпадает с~уравнением диффузии \cite{eleevbalk:bib:Leib}.

При этом на границе канала выполняются некоторые условия сопряжения.

Аналогичная ситуация имеет место для магнитной напряженности электромагнитного поля в~указанной выше неоднородной среде

\cite{eleevbalk:bib:TixSam}.

Большой интерес представляет изучение влияния вязкоупругих свойств нефти на различные технологические процессы ее добычи

\cite{eleevbalk:bib:Aziz, eleevbalk:bib:Ak}.

Если рассмотреть совместное движение различных несмешивающихся жидкостей в~трещинах и~пористых пластах с~учетом вязкоупругих

характеристик, то движение вязкоупругой и~вязкой жидкостей в~плоской горизонтальной трещине без учета поверхностных

явлений описывается одномерным гиперболическим уравнением и~уравнением теплопроводности с~интегро"=дифференциальными условиями

на границе раздела движущихся жидкостей.



В~монографии А.\,Г.~Шашкова \cite{eleevbalk:bib:Shash} строится структурная модель теплопроводности в~системе, составленной из

теплоизолированных с~боковой поверхности ограниченного и~полуограниченного стержней, имеющих одинаковую температуру.

На свободный конец системы поступает изменяющийся во времени тепловой поток.

Температурное поле в~ограниченном стержне описывается обычным уравнением теплопроводности, а~в~полуограниченном

"--- гиперболическим уравнением.

Теплофизические свойства стержней различны.

В~месте соприкосновения стержней имеет место идеальный тепловой поток.

Большой интерес представляет изучение математических моделей, описывающих влияние растительного покрова на теплообменные процессы в~почве и~приземном воздухе, при котором возникает необходимость исследования задачи для двух уравнений: уравнения Аллера переноса влаги, предполагающего бесконечную скорость распространения возмущения, и~уравнение А.\,В.~Лыкова, учитывающего конечную

его скорость.

При этом коэффициенты уравнений таковы, что они инвариантны относительно групп преобразований.



Достаточно полная библиография, посвященная постановке и~решению сопряженных задач для почв различных структур при различных

условиях на границе <<почва "--- воздух>>, имеется в монографии  А.\,Ф.~Чудновского \cite{eleevbalk:bib:Chudn}.

В~связи с изложенным возникает необходимость поиска корректно поставленных краевых задач, сформулированных одновременно для

волнового уравнения и~уравнения теплопроводности.

Среди ранних работ, тесно примыкающих по тематике к~данной статье, особо отметим работу Г.\,М.~Стручиной \cite{eleevbalk:bib:Str}.

Она посвящена задаче сопряжения двух уравнений по пространственной переменной.

Эта задача эквивалентно сводится к~интегральному уравнению Вольтерра второго рода, решение которого получено в~замкнутой форме.

Л.\,А.~Золина исследовала корректную постановку задачи сопряжения по временной переменной для таких уравнений и~изучила структурные

и~качественные свойства решения \cite{eleevbalk:bib:Zol}.

М.\,А.~Абдрахманов методом интегральных преобразований доказал однозначную разрешимость двумерной задачи сопряжения

\cite{eleevbalk:bib:Abdr}.

Е.\,А.~Островский рассмотрел также одномерную задачу сопряжения двух уравнений, когда в~граничные условия входят производные

по времени методом контурных интегралов \cite{eleevbalk:bib:Ostr}.

В.\,И.~Корзюк исследовал однозначную разрешимость задачи о~сопряжении уравнений гиперболического и~параболического типов

в~многомерном случае методами функционального анализа \cite{eleevbalk:bib:Kor}.



За последние несколько десятилетий в~математической литературе появилось значительное количество публикаций, посвященных задачам

сопряжения по временной переменной.

Достаточно полная библиография по этой теории содержится в~монографиях Т.\,Д.~Джураева \cite{eleevbalk:bib:Jur},

Т.\,Д.~Джураева, А.~Сопуева и М.~Мамажанова \cite{eleevbalk:bib:JurSop}, в~докторских диссертациях Д.~Базарова \cite{eleevbalk:bib:Baz},

В.\,А.~Елеева \cite{eleevbalk:bib:El}, О.\,А.~Репина \cite{eleevbalk:bib:Rep}, К.\,Б.~Сабитова \cite{eleevbalk:bib:Sab}.

В~приведенных выше работах задачи сопряжения двух уравнений по пространственной переменной в~основном изучались для

бесконечных или полубесконечных областей.



В~настоящей работе решаются задачи о~сопряжении гиперболического и~параболического уравнений по пространственной переменной

в~конечных областях.

Здесь отметим работу В.\,А.~Нахушевой \cite{eleevbalk:bib:Nah}, в~которой в~конечной области построена корректная математическая

модель задачи сопряжения,

найдены фундаментальные соотношения между температурой и~ее градиентом в~точке идеального контакта составной системы.



Рассматривается уравнение

\begin{equation} \label{el:1}

{0} = 

\begin{cases}

(-x)^m u_{tt}-a^2u_{xx}+b^2u, & 0 \leq m\leq2, \quad x<0, \\

u_{t}-c^2u_{xx}, & x>0

\end{cases}

\end{equation}

в конечной области $\Omega$, ограниченной отрезками $AB$, $BB_{0}$ и $A_{0}B_{0}$ прямых \linebreak $t=0$, $x=1$, $t=1$ соответственно

и~характеристиками $AC:\xi=t-(2/(m+2)\times(-x))^{(m+2)/2}=0$,  $A_0C:\eta=t+(2/(m+2)(-x))^{(m+2)/2}=1$, $\Omega_{1}=\Omega\bigcap(x>0)$,

$\Omega_{2}=\Omega\bigcap(x<0)$.



\begin{newthm}{Задача $\sc T_\alpha^\beta$}

Найти функцию $u(x,t),$ которая удовлетворяет следующим условиям\/{\rm:}

\begin{enumerate}

\item[\rm 1)] является регулярным решением уравнения \eqref{el:1} в области $\Omega$ при $x\neq 0$\/{\rm;}

\item[\rm 2)] удовлетворяет граничным условиям

\begin{equation} \label{el:2}

 u|_{AB}=\varphi_0(x), \quad  \Bigl[\alpha\frac{\partial u}{\partial x}+\beta

 u\Bigr]_{BB_0}=\varphi_1(t), \quad  u|_{AC}=\psi(x)

\end{equation}

и условиям сопряжения

\begin{equation} \label{el:3}

u(-0,t)=u(+0,t), \quad h_1\frac{\partial u}{\partial x}(-0,t)=h_2\frac{\partial u}{\partial x}(+0,t),

\end{equation}

где $\alpha,$ $\beta,$ $h_1,$ $h_2$ "--- константы$,$ $\varphi_0(x),$ $\varphi_1(t),$ $\psi(x)$ "--- заданные функции\rm.

\end{enumerate}

\end{newthm}



Справедлива

\begin{theorem}[1]Пусть выполняются условия\/{\rm:}

\begin{enumerate}

\item[\rm 1)] $0<m<2,$ $\alpha=0,$ $\beta=1$\/{\rm;}

$\varphi_0(0)=\psi(0),$ $\varphi_0(x)\in C^1[0, 1],$ $\varphi_1(t)\in C[0, 1]$\/{\rm;} 

$\psi(x)\in C^4[-1/2, 0],$ $a=1,$ $c=1,$ $b=0,$ $\varphi_0(0)=\varphi_1(0)=0,$ $\varphi_1(1)=\varphi_1(0)$\/{\rm;}

\item[\rm 2)] $\alpha=0,$ $\beta=1;$ $\varphi_0(x)\in C^2[0, 1],$ $\varphi_1(t)\in C[0, 1],$ $\psi(x)\in C^4[-1/2, 0],$

$\varphi_0(0)=\varphi'_0(0)=\psi(0)=0m$ $\varphi_0(1)=\varphi_1(0)$\/{\rm;} $a=1,$ $c=1,$ $b=0,$ $0<m<2.$

\end{enumerate}

Тогда задача $T_0^1$ имеет$,$ и притом единственное$,$ решение\rm.

\end{theorem}



\begin{proof}

Рассмотрим условие 1). Переходя к решению поставленной задачи, воспользуемся первым условием сопряжения из (\ref{el:3}) и~обозначим через $\mu(t)$

значение функции $u(x,t)$ на прямой $x=0$:

\begin{equation} \label{el:4}

\mu(t)=u(-0,t)=u(+0,t).

\end{equation}



Решение первой краевой задачи \eqref{el:1}, \eqref{el:2}, \eqref{el:4} в~$\Omega_1$ имеет вид

\begin{multline} \label{el:5}

u(x,t)=\int_0^1\varphi_0(\xi)G(\xi,0;x,t)d\xi+\int_0^t\mu(\eta)G_\xi(0,\eta;x,t)d\eta-

\\

-\int_0^t\varphi_1(\eta)G_\xi(1,\eta;x,t)d\eta,

\end{multline}

где

\begin{multline} \label{el:6}

G(\xi,\eta;x,t)=\frac{1}{2\sqrt{\pi(t-\eta)}}\sum\limits_{n=-\infty}^{n=\infty}\Bigl\{\exp\Bigl[-\frac{(x-\xi+2n)^2}{4(t-\eta)}\Bigr]-

\\

-\exp\Bigl[-\frac{(x+\xi+2n)^2}{4(t-\eta)}\Bigr]\Bigr\}

\end{multline}

"--- функция Грина первой краевой задачи для уравнения (\ref{el:1}) при $x>0$.

В~силу свойств функции Грина (\ref{el:6}) из (\ref{el:5}) получаем \cite{eleevbalk:bib:Eleev}

\begin{equation}\label{el:7}

\nu^+(t)=-\frac{1}{\sqrt{\pi}}\int_0^t\frac{\mu'(\eta)d\eta}{(t-\eta)^{1/2}}+\int_0^tk_0(t,\eta)

\mu'(\eta)d\eta+\psi_0(t),

\end{equation}

где

\begin{equation} \label{el:8}

k_0(t,\eta)=-\frac{1}{\sqrt{\pi}}\mathop{\sum{\rule{0mm}{10pt}}'}\limits_{n=-\infty}^{n=\infty} (t-\eta)^{-1/2}\exp\Bigl(-\frac{n^2}{t-\eta}\Bigr),

\end{equation}



$$\psi_0(t)=\int_0^1\varphi_0(\xi)G_x(\xi,0;0,t)d\xi-\int_0^t\varphi_1(\eta)G_{\xi x}(1,\eta;0,t)dt,$$

где $\sum'$ означает, что $n\neq 0$, $\nu^+(t)=\lim\limits_{x\rightarrow +0}u_x(x,t)$.



Учитывая свойства заданных функций $\varphi_0(x)$ и~$\varphi_1(t)$, можем заключить, что $\psi_0(t)\in C^1[0, 1]\bigcap C^2]0, 1[$.

Найдём теперь функциональное соотношение между $\mu(t)$ и~$\nu^-(t)$, которое приносится из гиперболической части $\Omega_2$

на линию $x=0$. Для этого в~уравнении (\ref{el:1}) при $x<0$ перейдём к характеристическим координатам:

$$\xi=t-(2/(m+2))(-x)^{(m+2)/2}, \quad \eta=t+(2/(m+2))(-x)^{(m+2)/2}.$$



В результате уравнение (\ref{el:1}) в~области его гиперболичности перейдёт в~уравнение Эйлера"--~Дарбу"--~Пуассона:

\begin{equation} \label{el:9}

u_{\xi\eta}+\frac{\bar{\beta}}{\eta-\xi}(u_\xi-u_\eta)=0, \quad 0<\xi<\eta<1,

\end{equation}

где $2\bar{\beta}=m/(m+2)$, а краевое условие $u|_{AC}=\psi(x)$ "--- в~условие

\begin{equation} \label{el:10}

u|_{\xi=0}=\psi\Bigl[-\Bigl(\frac{m+2}{4}\Bigr)^{1-2\bar{\beta}}\eta^{1-2\bar{\beta}}\Bigr]=\varphi(\eta), \quad 0\leq\eta\leq1.

\end{equation}

При этом

\begin{equation} \label{el:11}

\Big(\frac{m+2}{2}\Bigr)^{2\bar{\beta}}\left[(\eta-\xi)^{2\bar{\beta}}(u_\xi-u_\eta)\right]_{\eta=\xi}=\nu^-(\xi).

\end{equation}



Решение задачи Дарбу \eqref{el:9}--\eqref{el:11} выражается формулой

\begin{multline} \label{el:12}

u(\xi,\eta)=\gamma_0\int_0^\xi(\xi-z)^{-\bar{\beta}}(\eta-z)^{-\bar{\beta}}\nu^-(z)dz+

\\

+\int_0^\eta[\varphi'(z)+\bar{\beta}z^{-1}\varphi(z)]H(0,z;\xi,\eta)dz,

\end{multline}

где $\gamma_0=(-2^{2\bar{\beta}-1}\rho)/(m+2)^{2\bar{\beta}}$, $\rho=\Gamma(\bar{\beta})/[\Gamma(1-\bar{\beta})\Gamma(2\bar{\beta})]$, $H(0,z;\xi,\eta)=\rho

z^{2\bar{\beta}}\eta^{-\bar{\beta}}\times(\xi-z)^{-\bar{\beta}}$ "--- след функции Грина"--~Адамара уравнения (\ref{el:9}) \cite{eleevbalk:bib:Eleev}.



В равенстве (\ref{el:12}) положим $\eta=\xi=t$, $0<t<1$.

В~результате получим функциональное соотношение между $\mu(t)$ и $\nu^-(t)$:

\begin{equation} \label{el:13}

\mu(t)=\gamma_0\int_0^t(t-\eta)^{-2\bar{\beta}}\nu^-(\eta)d\eta+\psi_1(t),

\end{equation}

где

\begin{equation} \notag

\psi_1(t)=\rho t^{-\beta}\int_0^t\eta^{2\bar{\beta}}(t-\eta)^{-\bar{\beta}}[\varphi'(\eta)+\bar{\beta}\eta^{-1}\varphi(\eta)]d\eta,

\end{equation}

причём $\alpha^j\varphi(\eta)/d\eta^j=\eta^{1-2\bar{\beta}}\varphi^{*}_j(\eta)$, где

$\varphi^{*}_j(\eta)\in C[0, 1]$, $j$ принимает значения $0, 1, 2, 3$.



С учётом этого легко заметить, что $\psi_1(t) \in C[0, 1]\bigcap C^3]0, 1[$, $\psi_1(t)=t^{1-2\bar{\beta}} {\sf O}(1)$, где символ ${\sf O}(1)$

означает ограниченную величину.

Перепишем уравнение (\ref{el:13}) в виде

\begin{equation} \label{el:14}

D_{0t}^{2\bar{\beta}-1}\nu^-(t)=\bar{\gamma_0}\mu(t)+\bar{\gamma_0}\psi_1(t),

\end{equation}

где $\bar{\gamma_0}=1/[\gamma_0\Gamma(1-2\bar{\beta})]$, $D_{0t}^l$ "--- оператор дробного интегрирования порядка ${-l}$ при $l<0$

и~обобщенная в~смысле Лиувилля производная порядка $l$ при $l>0$.



К обеим частям равенства (\ref{el:14}) применим оператор $D_{0t}^{1-2\bar{\beta}}$, обратный оператору $D_{0t}^{2\bar{\beta}-1}$.

Будем иметь

\begin{equation} \label{el:15}

\nu^-(t)=\bar{\gamma_0}D_{0t}^{1-2\bar{\beta}}\mu(t)+\bar{\gamma_0}D_{0t}^{1-2\bar{\beta}}\psi_1(t).

\end{equation}



В силу второго условия сопряжения из (\ref{el:3}) получим

\begin{equation} \label{el:16}

D_{0t}^{-2\bar{\beta}}\mu'(t)+\lambda D_{0t}^{-1/2}\mu'(t)=\frac{1}{\Gamma(2\bar{\beta})}\biggl[\overline{\psi}_0(t)+\int_0^tk_0(t,\eta)\mu'(\eta)d\eta\biggr],

\end{equation}

где $\lambda=-[\sqrt{\pi}\Gamma(1-2\bar{\beta})\bar{\gamma}_0h_2]/\sqrt{\pi}h_1$,

$\bar{\psi}_0(t)=\psi_0(t)-\frac{h_1}{\gamma_0\Gamma(1-2\bar{\beta})}D_{0t}^{1-2\bar{\beta}}\bar{\psi}_1(t)$.



Действуя на обе части равенства (\ref{el:16}) оператором $D_{0t}^{2\bar{\beta}}$, получим

\begin{equation} \label{el:17}

\mu'(t)+\bar{\lambda}\int_0^t\frac{\mu'(\eta)}{(t-\eta)^{2\bar{\beta}+1/2}}d\eta=\frac{1}{\Gamma(2\bar{\beta})}D_{0t}^{2\bar{\beta}}\biggl[\overline{\psi}_0(t)+\int_0^tk_0(t,\eta)\mu'(\eta)d\eta\biggr],

\end{equation}

где $\bar{\lambda}={\lambda}/\Gamma(1/2-2\bar{\beta})$.



Резольвента ядра $(t-\eta)^{-2\bar{\beta}-1/2}$ оператора Вольтерра в левой части (\ref{el:17}) имеет вид \cite{eleevbalk:bib:Mis}

\begin{equation} \label{el:18}

R(t-\eta,\bar{\lambda})=\sum\limits_{n=1}^\infty

\bar{\lambda}^{n-1}\frac{[\Gamma(1/2-2\bar{\beta})]^{n}}{\Gamma[n(1/2-2\bar{\beta})]}(t-\eta)^{\eta(1/2-2\bar{\beta})-1}.

\end{equation}



После обращения интегрального оператора Вольтерра уравнение (\ref{el:17}) примет вид

\begin{multline} \label{el:19}

\mu'(t)+\bar{\bar{\lambda}}\int_0^tR(t-\eta,\bar{\lambda})D_{0t}^{2\bar{\beta}}\int_0^tk_0(t,\eta_1)\mu'(\eta_1)d\eta_1d\eta=

\\

=\bar{\bar{\lambda}}\int_0^tR(t-\eta,\bar{\lambda})D_{0t}^{2\bar{\beta}}\bar{\psi}_0(t)d\eta,

\end{multline}

где $\bar{\bar{\lambda}}=\bar{\lambda}/\Gamma(2\bar{\beta})$.



Принимая во внимание, что в равенстве (\ref{el:19})

\begin{multline*}

D_{0t}^{2\bar{\beta}}\int_0^tk_0(t,\eta_1)\mu'(\eta_1)d\eta_1=

\\

=\frac{1}{\Gamma(1-2\bar{\beta})}\frac{d}{dt}\int_0^t(t-\eta_1)^{3/2-2\bar{\beta}} k_1(t,\eta_1)\mu'(\eta_1)d\eta_1=

\\

=\frac{1}{\Gamma(1-2\bar{\beta})}\int_0^t[(3/2-2\bar{\beta})(t-\eta_1)^{1/2-2\bar

{\beta}}k_1(t,\eta_1)+

\\

+(t-\eta_1)^{-1/2-2\bar{\beta}}k_2(t,\eta_1)]\mu'(\eta_1)d\eta_1=\frac{1}{\Gamma(1-2\beta)}\int_0^t\frac{N(t,\eta_1)\mu'

(\eta_1)d\eta_1}{(t-\eta_1)^{2\bar{\beta}+1/2}},

 \end{multline*}

где

$$k_1(t,\eta_1)=\int_0^1z^{1/2}(1-z)^{-2\bar{\beta}}\mathop{\sum{\rule{0mm}{10pt}'}} \limits_{n=-\infty}^{n=\infty}\exp\Bigl(-\frac{n^2}{(t-\eta_1)z}\Bigr)dz,$$

$$k_2(t,\eta_1)=n^2\int_0^1z^{-1/2}(1-z)^{-2\bar{\beta}}\mathop{\sum{\rule{0mm}{10pt}'}} \limits_{n=-\infty}^{n=\infty}\exp\Bigl(-\frac{n^2}{(t-\eta_1)z}\Bigr)dz,$$

\begin{multline*}N(t,\eta_1)=(3/2-2\bar{\beta})(t-\eta_1)k_1(t,\eta_1)+k_2(t,\eta_1)\in C[0, 1]\times

\\

\times[0, 1]\bigcap C^2]0, 1[{}\times{}]0, 1[{}, \;

D_{0t}^{2\bar{\beta}}\bar{\psi}_0(t)\in C[0, 1]\bigcap C^2]0, 1[{},

\end{multline*}

получим интегральное уравнение Вольтерра второго рода с

интегрируемой особенностью:

\begin{equation} \label{el:20}

\mu'(t)+\bar{\bar{\lambda}}\int_0^t\frac{\bar{N}(t,\eta)\mu'(\eta)}{(t-\eta)^{2\bar{\beta}-1/2}}d\eta=q(t),

\end{equation}

где

$$ \bar{N}(t,\eta)=\int_0^1\frac{R\bigl((t-\eta)(1-z),\bar{\lambda}\bigr)N(\eta+(t-\eta)z,\eta)}{z^{2\bar{\beta}+1/2}}dz,$$

$$ q(t)=\bar{\bar{\lambda}}\int_0^tR(t-\eta,\bar{\lambda})D_{0t}^{2\bar{\beta}}\bar{\psi}_0(t)d\eta.$$



В силу условий, наложенных на функцию $\psi(x)$, можем заключить, что $q(t)\in C[0, 1]\bigcap C^2]0, 1[$, следовательно, уравнение

(\ref{el:20}) однозначно и безусловно разрешимо в классе $C[0, 1]$.



Обозначив через $T(t,\eta)$ резольвенту уравнения (\ref{el:20}) и~обратив его как интегральное уравнение Вольтерра второго рода, 

получим

$$\mu'(t)=q(t)+\int_0^tT(t,\eta)q(\eta)d\eta$$

или

\pagebreak



$$\mu(t)=\int_0^t\biggl[1+\int_{\eta}^tT(t,\eta_1)d\eta_1\biggr]q(\eta)d\eta.$$



После определения функции $\mu(t)$ по формуле (\ref{el:7}) находим $\nu^+(t)$, а из равенства (\ref{el:15}) "---  $\nu^-(t)$.

Таким образом, решение задачи $T_0^1$ в области $\Omega_1$ даётся формулой (\ref{el:5}), а~в~$\Omega_2$ "--- формулой (\ref{el:12}).



Однозначная разрешимость случая 2 доказывается аналогично предыдущему случаю с использованием свойств функции Грина смешанной краевой

задачи (\ref{el:1}) при $x>0$: $u(x, 0)=\varphi_0(x)$, $u(1,t)=\varphi_1(t)$, $u_x(0,t)=\varphi_2(t)$.

\end{proof}



\smallskip



\begin{theorem}[2]Пусть выполняются условия

\begin{enumerate}

\item[\rm 1)] $\alpha=1${\rm,} $\beta=0${\rm,} $\varphi_0(x)\in C^1[0, 1]${\rm,} $\varphi_1(t)\in C[0, 1]${\rm,} $\psi(x)\in C^4[-1/2, 0]${\rm,}

$\varphi_0(0)=\psi(0)=0${\rm,} $\varphi'_0(1)=\varphi_1(0)${\rm,} $0<m<2$ 

\end{enumerate}

или

\begin{enumerate}

\item[\rm 2)] $\alpha=1${\rm,} $\beta=0${\rm;} $\varphi_0(x)\in C^2[0, 1]${\rm,} $\varphi_1(t)\in C[0, 1]${\rm,} $\varphi_0(0)=\varphi'_0(0)=\psi(0)=0${\rm,}

$\varphi_0(1)=\varphi_1(0)${\rm,} $0<m<2$\rm.

\end{enumerate}

Тогда задача $T_1^0$ имеет$,$ и~притом единственное$,$ решение\rm.

\end{theorem}



\begin{proof}

Предположим, что решение задачи $T_1^0$ существует.

Тогда в силу дифференциальных свойств функции Грина \cite{eleevbalk:bib:Eleev}

\begin{multline*}

G(\xi,\eta;x,y)=\frac{\pi^{-1/2}}{2}(t-\eta)^{-1/2}\sum\limits_{n=-\infty}^{n=\infty}\Bigl\{\exp\Bigl[-\frac{(x-\xi+2n)^2}{4(t-\eta)}\Bigr]-\\

-\exp\Bigl[-\frac{(x+\xi+2n)^2}{4(t-\eta)}\Bigr]+\exp\Bigl[-\frac{(x+\xi-2+2n)^2}{4(t-\eta)}\Bigr]-\exp\Bigl[-\frac{(x-\xi-2+2n)^2}{4(t-\eta)}\Bigr]\Bigr\}

\end{multline*}

смешанной краевой задачи $u(x, 0)=\varphi_0(x)$, $u(0,t)=\mu(t)$, $u_x(1,t)=\varphi_1(t)$ уравнения (\ref{el:1}) при $x>0$

функция $u(x,t)$ является решением нагруженного интегрального уравнения

\begin{multline}\label{el:21}

u(x,t)=\int_0^1\varphi_0(\xi)G(\xi,0;x,t)d\xi+\int_0^tu(0,\eta)G_\xi(0,\eta;x,t)d\eta+\\

+\int_0^1\varphi_1(\eta)G(1,\eta;x,t)d\eta.

\end{multline}

Функциональное соотношение между $\mu(t)$ и $\nu^+(t)$ на линии $x=0$ имеет вид

\begin{equation} \label{el:22}

\nu^+(t)=-D_{0t}^{1/2}\mu(t)+\int_0^tk_3(t,\eta)\mu(t)dt+\psi_2(t),

\end{equation}

где

\begin{multline}\label{el:23}

k_3(t,\eta)=\frac{1}{2\sqrt{\pi}}(t-\eta)^{-5/2}\biggl\{(2-t+\eta)\exp[-1/(t-\eta)]+\\

+\mathop{\sum\rule{0mm}{11pt}{'}}\limits_{n=-\infty}^{n=\infty}\bigl((t-\eta-2n^2)\exp[-n^2/(t-\eta)]+\\

+(2(n-1)^2-t+\eta)\exp[-(n-1)^2/(t-\eta)]\bigr)\biggr\},

\end{multline}



\pagebreak



\begin{equation} \label{el:24}

\psi_2(t)=\int_0^1\varphi_0(\xi)G_\xi(\xi, 0;0,t)d\xi+\int_0^t\varphi_1(\eta)G_x(1,\eta;x,t)d\eta.

\end{equation}



В силу второго условия сопряжения из (\ref{el:3}) и равенств (\ref{el:15}), (\ref{el:22}) получим интегральное уравнение для

определения $\mu(t)$:

$$ \mu(t)+\lambda D_{ot}^{2\bar{\beta}-1/2}\mu(t)+\lambda\int_0^t\bar{k_3}(t,\eta)\mu(t)dt=\lambda

D_{ot}^{2\bar{\beta}-1}\psi_2(t)-\psi_1(t),$$

где $\lambda=h_2/(h_1\bar{\gamma_0})$,

$$\bar{k}_3(t,\eta)=\frac{1}{\Gamma(1-2\bar{\beta})}\int_{\eta}^t(t-\eta_1)^{-2\bar{\beta}}k_3(\eta_1,\eta)d\eta_1.$$



Из представлений (\ref{el:23}), (\ref{el:24}) функций $k_3(t,\eta)$, $\psi_2(t)$ легко заключить, что

$\bar{k}_3(t,\eta)\in C[0, 1]\times [0, 1]$, $\psi_2(t)\in C[0, 1]$.



Поступая таким же образом, как для уравнения (\ref{el:20}), находим единственным образом функцию $\mu(t)$, а затем из равенств

(\ref{el:22}) и~(\ref{el:15}) "--- функции $\nu^+(t)$ и $\nu^-(t)$.



Следовательно, решение задачи $T_1^0$ в области $\Omega_1$ определяется формулой (\ref{el:21}), а~в~$\Omega_2$ "--- формулой 

(\ref{el:12}). Случай 2 доказывается аналогично.

\end{proof}



\smallskip



\begin{theorem}[3]Пусть выполнены следующие условия{\rm:} $\alpha=1,$ $\beta=-\gamma={\rm const}>0,$ $2\bar{\beta}=1/2,$ т.~е. $m=2,$ $u(x,0)=0,$ $x\geq0,$ $\varphi_1(t)\in C[0, 1],$ $\varphi'_0(1)-\gamma\varphi_0(1)=\varphi_1(0).$ Тогда задача

$T_1^{-\gamma}$ имеет$,$ и притом единственное$,$ решение\rm.

\end{theorem}



\begin{proof}

Решение задачи $T_1^{-\gamma}$ в области $\Omega_1$ будем искать в~виде суммы тепловых потенциалов двойного слоя и простого слоя:

\begin{equation} \label{el:25}

u(x,t)=\int_0^t\theta(\eta)\frac{x}{2\sqrt{\pi(t-\eta)}}e^{-\frac{x^2}{4(t-\eta)}}\frac{d\eta}{2(t-\eta)}+\int_0^t\nu^+(\eta)\frac{e^{-\frac{(x-1)^2}{4(t-\eta)}}d\eta}{2\sqrt{\pi(t-\eta)}}.

\end{equation}

Удовлетворяя (\ref{el:25}) второму граничному условию из (\ref{el:2}) и  условию (\ref{el:4}),

получим систему из интегральных уравнений Вольтерра первого и второго рода:

\begin{equation} \label{el:26}

\begin{array}{l}

\displaystyle

\mu(t)=\frac{\theta(t)}{2}+\frac{1}{2}\int_0^t\nu^+(t)\frac{1}{\sqrt{\pi(t-\eta)}}e^{-\frac{1}{4(t-\eta)}}d\eta,

\\[2mm] 

\displaystyle

\varphi_1(t)=\frac{\nu^+(t)}{2}-\frac{\gamma}{2}\int_0^t\frac{\nu^+(\eta)d\eta}{\sqrt{\pi(t-\eta)}}-\frac{2\gamma+1}{4}\int_0^t\theta(\eta)e^{-\frac{1}{4(t-\eta)}}\frac{d\eta}{t-\eta}.

\end{array}

\end{equation}

Соотношение между $\mu(t)$ и $\nu^-(t)$, приносимое на линию $x=0$, имеет вид

\begin{equation} \label{el:27}

\mu(t)=\gamma_0\int_0^t\frac{\nu^-(\eta)d\eta}{(t-\eta)^{1/2}}+\psi_1(t).

\end{equation}



Перепишем второе уравнение системы (\ref{el:26}) в виде

\begin{equation} \label{el:28}

\nu^+(t)=\lambda\int_0^t\frac{\nu^+(\eta)}{(t-\eta)^{1/2}}d\eta+F(t),

\end{equation}

где

$$F(t)=2\varphi_1(t)+\frac{2\gamma+1}{2}\int_0^t\theta(\eta)\frac{1}{t-\eta}e^{-\frac{1}{4(t-\eta)}}d\eta,

\quad \lambda=\gamma/\sqrt{\pi}.$$



Элементарные вычисления показывают, что в данном случае повторные

ядра уравнения (\ref{el:28}) выражаются так \cite{eleevbalk:bib:Str}:

$K_n(t-\eta)=\pi^{\frac{n}{2}}(t-\eta)^{\frac{n-2}{2}}/\Gamma(n/2)$.



Тогда резольвента ядра интегрального уравнения (\ref{el:28}) имеет вид



\begin{equation} \label{el:29}

R(t-\eta)=\sum\limits_{n=1}^\infty\Bigl[\lambda^n\pi^{\frac{n}{2}}(t-\eta)^{\frac{n-2}{2}}/\Gamma(n/2)\Bigr].

\end{equation}



Решение уравнения (\ref{el:28}) даётся формулой

\begin{equation} \label{el:30}

\nu^+(t)=\int_0^tR(t-\eta)F(\eta)d\eta+F(t).

\end{equation}



Учитывая в равенстве (\ref{el:30}) значение функции $F(t)$ и выполняя несложные

преобразования, получим

\begin{equation} \label{el:31}

\nu^+(t)=\int_0^tZ(t,\eta)\theta(\eta)d\eta+\bar{\varphi}_1(t),

\end{equation}

где

$$Z(t,\eta)=\int_0^1R[(t-\eta)(1-z)]\frac{1}{z}e^{-\frac{1}{4(t-\eta)z}}dz+\frac{1}{t-\eta}e^{-\frac{1}{4(t-\eta)}},$$

$$\bar{\varphi}_1(t)=2\biggl(\varphi_1(t)+\int_0^tR(t-\eta)\varphi_1(\eta)d\eta\biggr).$$



Подставляя (\ref{el:27}) в первое уравнение системы (\ref{el:26}) и учитывая второе условие сопряжения из (\ref{el:3}), получим

\begin{equation} \label{el:32}

\int_0^t\frac{G(t,\eta)}{\sqrt{t-\eta}}\nu^+(\eta)d\eta=\frac{\theta(t)}{2}-\psi_1(t),

\end{equation}

где

$$G(t,\eta)=\dfrac{h_2\gamma_0}{h_1}-\dfrac{1}{2\sqrt{\pi}}e^{-\frac{1}{4(t-\eta)}},$$ 

причём $G(t,t)=(h_2\gamma_0)/h_1$.



Следовательно, уравнение (\ref{el:32}) можно свести к интегральному уравнению Вольтерра второго рода:

\begin{equation} \label{el:33}

\nu^+(t)+\lambda_0\int_0^t\frac{\partial N(t,\eta)}{\partial

t}\nu^+(\eta)d\eta=\lambda_1\int_0^t\frac{\theta'(\eta)d\eta}{\sqrt{t-\eta}}+\lambda_2\int_0^t\frac{\psi'_1(\eta)d\eta}{\sqrt{t-\eta}},

\end{equation}

где

$$N(t,\eta)=\int_0^1\frac{G[\eta+(t-\eta)z,\eta]dz}{z^{1/2}(1-z)^{1/2}},$$

$$\lambda_0=h_1/(\gamma_0h_2), \quad \lambda_1=(2h_1)/(\gamma_0h_2), \quad \lambda_2=h_1/h_2.$$

Исключая из системы (\ref{el:31}), (\ref{el:33}) $\nu^+(t)$ и выполняя несложные преобразования, получим

\begin{equation} \label{el:34}

\int_0^t\frac{\theta'(\eta)d\eta}{\sqrt{t-\eta}}=\frac{1}{\lambda_1}\int_0^t\bar{Z}(t,\eta)\theta(\eta)d\eta+\bar{\bar{\varphi}}_1(t),

\end{equation}

где

$$\bar{Z}(t,\eta)=Z(t,\eta)+\lambda_0\int_{\eta}^t\frac{\partial

N(t,\eta_1)}{\partial t}Z(\eta_1,\eta)d\eta_1,$$

$$\bar{\bar{\varphi}}_1(t)=\bar{\varphi}_1(t)+\lambda_0\int_0^t\frac{\partial

k_1(t,\eta)}{\partial t}\bar{\varphi}_1(\eta)d\eta,$$

причём $\bar{Z}(t,t)=0$. Обращая (\ref{el:34}) как интегральное уравнение Абеля относительно $\theta'(t)$, а затем интегрируя обе  части полученного равенства от $0$ до $t$, будем иметь

\begin{equation} \label{el:35}

\theta(t)-\bar{\lambda}_1\int_0^tM(t,\eta)\theta(\eta)d\eta=r(t),

\end{equation}

где

$$M(t,\eta)=\int_{\eta}^t\frac{\bar{Z}(\eta_1,\eta)d\eta_1}{\sqrt{t-\eta_1}}, \quad r(t)=-\frac{1}{\pi}\int_0^t\frac{\bar{\bar{\varphi}}_1(\eta)d\eta}{\sqrt{t-\eta}}, \quad

\bar{\lambda}_1=1/(\lambda_1\pi).$$



В силу представлений $N(t,\eta)$, $Z(t,\eta)$, а следовательно $\bar{Z}(t,\eta)$ и свойств функции $\varphi_1(t)$, нетрудно

проверить, что $M(t,\eta)\in C[0,1]\times[0,1]$, $r(t)\in C[0,1]$.



Обращая интегральный оператор Вольтерра через резольвенту $T(t,\eta)$, получим

\begin{equation} \label{el:36}

\theta(t)=r(t)+\bar{\lambda}_1\int_0^tT(t,\eta)r(\eta)d\eta.

\end{equation}

Подставляя значение $\theta(t)$ из соотношения (\ref{el:36}) в равенство (\ref{el:31}), находим $\nu^{+}(t)$, затем из первого 

равенства системы (\ref{el:26}) определяем $\mu(t)$, а из второго условия сопряжения (\ref{el:3}) получаем $\nu^-(t)$.



Следовательно, решение задачи $T_\alpha^{-\gamma}$ в области $\Omega_1$ даётся формулой (\ref{el:25}), а в $\Omega_2$ "--- равенством  (\ref{el:12}).

\end{proof}



\smallskip



\begin{theorem}[4]Пусть выполнены условия {\rm1} теоремы {\rm1:} $a=1,$ $c=1,$ $b\neq 0,$ $m=0.$ Тогда задача $T_0^1$ имеет$,$ и притом единственное$,$ решение\rm.

\end{theorem}



\begin{proof}

Решение задачи Коши $u(0,t)=\mu(t)$, $u_x(0,t)=\nu^{-}(t)$ в области

$\Omega_2$ даётся формулой

\begin{multline} \label{el:37}

u(x,t)=\frac{\mu(x+t)+\mu(t-x)}{2}+\frac{1}{2}\int_{t-x}^{t+x}\nu^-(\eta)J_0(b\sqrt{x^2-(t-\eta)^2})d\eta+\\

+\frac{bx}{2}\int_{t-x}^{t+x}\frac{J_1(b\sqrt{x^2-(t-\eta)^2})}{\sqrt{x^2-(t-\eta)^2}}\mu(\eta)d\eta,

\end{multline}

где $J_0(z)$ и $J_1(z)$ "--- модифицированные функции Бесселя первого рода. Удовлетворяя (\ref{el:37}) третьему условию из 

(\ref{el:2}), получаем функциональное соотношение между $\mu(t)$ и $\nu^-(t)$, приносимое из области $\Omega_2$ на линию $x=0$:

$$

\mu(t)+\int_0^t\frac{t}{\eta}\frac{\partial}{\partial t}J_0\bigl(b\sqrt{\eta(t-\eta)}\bigr)\mu(\eta)d\eta=2\psi(t/2)+\int_0^t\nu^-(\eta)J_0\bigl(b\sqrt{\eta(t-\eta)}\bigr)d\eta.

$$



Отсюда после обращения относительно $\mu(t)$ и ряда преобразований получим \cite{eleevbalk:bib:Salah}

\begin{equation} \label{el:38}

\mu(t)=\int_0^tJ_0[b(t-\eta)]\nu^-(\eta)d\eta+\phi(t),

\end{equation}

где

$$\phi(t)=2\biggl(\psi(\frac{t}{2})-\int_0^t\psi\Bigl(\frac{\eta}{2}\Bigr)\frac{\partial}{\partial \eta}J_0\bigl(b\sqrt{t(\eta-t)}\bigr)d\eta\biggr),$$

$J_0(z)$ "--- функция Бесселя первого рода нулевого порядка. Обращение интегрального уравнения (\ref{el:38}) относительно 

$\nu^-(t)$ имеет вид

\begin{equation} \label{el:39}

\nu^-(t)=\mu'(t)+b\int_0^t\frac{J_1[b(t-\eta)]}{t-\eta}\mu(\eta)d\eta+\phi_1(t),

\end{equation}

где

$$\phi_1(t)=\phi'(t)+b\int_0^t\frac{J_1[b(t-\eta)]}{t-\eta}\phi(\eta)d\eta,$$

$J_1(z)$ "--- функция Бесселя первого рода первого порядка. В силу второго условия сопряжения из (\ref{el:3}) получим

\begin{equation} \label{el:40}

\mu'(t)+\lambda\int_0^t\frac{\mu'(\eta)d\eta}{(t-\eta)^{1/2}}=\bar{\phi}_1(t)+\int_0^t\bar{T}(t,\eta)\mu(\eta)d\eta,

\end{equation}

где

$$\bar{\phi}_1(t)=\frac{h_2}{h_1}\psi_0(t)-\phi_1(t), \quad \bar{T}(t,\eta)=-\frac{h_2}{h_1}\frac{\partial k_0(t,\eta)}{\partial

\eta}-c\frac{J_1[b(t-\eta)]}{t-\eta}.

$$



С помощью резольвенты (\ref{el:29}) оператора в левой части равенства (\ref{el:40}) можем записать

\begin{equation} \label{el:41}

\mu'(t)=\bar{\bar{\phi}}_1(t)+\int_0^t\bar{\bar{T}}(t,\eta)\mu(\eta)d\eta,

\end{equation}

где

$$\bar{\bar{\phi}}_1(t)=\bar{\phi}_1(t)+\int_0^tR(t-\eta)\bar{\phi}_1(\eta)d\eta, \quad

\bar{\bar{T}}(t,\eta)=\bar{T}(t,\eta)+\int_{\eta}^tR(t-\eta)\bar{T}(\eta_1,\eta)d\eta_1.$$



Путём интегрирования в пределах от $0$ до $t$ равенства (\ref{el:41}) после некоторых преобразований приходим к интегральному 

уравнению Вольтерра второго рода:

\begin{equation} \label{el:42}

\mu(t)+\int_0^tP(t,\eta)\mu(\eta)d\eta=n(t),

\end{equation}

где

$$P(t,\eta)=\int_{\eta}^t\bar{\bar{T}}(\eta_1,\eta)d\eta_1, \quad n(t)=\int_0^t\bar{\bar{\phi}}_1(\eta)d\eta.$$



В силу свойств функций $P(t,\eta)$, $n(t)$ заключаем, что существует единственное решение интегрального уравнения (\ref{el:42}) 

в~классе $C[0,1]$.



После определения $\mu(t)$ по формуле (\ref{el:39}) находим $\nu^-(t)$. Следовательно, решение задачи $T_0^1$ в областях 

$\Omega_1$ и $\Omega_2$ находится по формулам (\ref{el:37}) и (\ref{el:6}) соответственно.

\end{proof}



\begin{newthm}{Задача 1 (N)} Задача N отличается от задачи $T_\alpha^\beta$ тем, что второе условие сопряжения из \eqref{el:3}

заменяется условием

\begin{equation} \label{el:43}

\lambda_1\frac{\partial u}{\partial

x}(-0,t)=\frac{\lambda_2}{\tau_q}\int_0^t\frac{\partial

u}{\partial x}(x,\eta)\Big|_{x=0}\exp\Bigl(-\frac{t-\eta}{\tau_q}\Bigr)d\eta,

\end{equation}

где $\tau_q$ "--- время релаксации теплового потока.

\end{newthm}



Учитывая (\ref{el:7}), (\ref{el:39}) в равенстве (\ref{el:43}) и выполняя несложные

преобразования, получим

\begin{equation} \label{el:44}

\mu'(t)+\lambda\int_0^t\Pi(t,\eta)\mu'(\eta)d\eta=e(t)-b\int_0^t\frac{J_1[b(t-\eta)]}{t-\eta}\mu(\eta)d\eta,

\end{equation}

где

$$\Pi(t,\eta)=\int_{\eta}^t\bigl[(\eta_1-\eta)^{-1/2}+k_0(\eta_1,\eta)\bigr]\exp\Bigl(-\frac{t-\eta_1}{\tau_q}d\eta_1\Bigr),$$

$$e(t)=\lambda\int_0^t\psi_0(\eta)\exp\Bigl(-\frac{t-\eta}{\tau_q}\Bigr)d\eta-\psi(t), \quad \lambda=\lambda_2/(\lambda_1\sqrt{\pi}\tau_q).$$



В силу свойств заданных функций $\varphi_0(x)$, $\varphi_1(t)$ и $\psi(x)$ заключаем, что $\Pi(t,\eta)\in C[0, 1]\bigcap C^2]0, 1[,

e(t)\in C[0, 1]\bigcap C^2]0, 1[$ по переменным $t$ и $\eta$. Если правую часть уравнения (\ref{el:44}) считать известной и 

обозначить через $N(t,\eta)$ резольвенту ядра $\Pi(t,\eta)$, то оно может быть представлено в виде

\begin{equation} \label{el:45}

\mu'(t)=\rho(t)+\lambda\int_0^t\bar{N}(t,\eta)\mu(\eta)d\eta,

\end{equation}

где

\begin{multline*}

\rho(t)=\lambda\int_{0}^t\psi_0(\eta)\exp\Bigl(-\frac{t-\eta}{\tau_q}\Bigr)d\eta-\psi_1(t)+\\

+\lambda\int_{0}^tN(t,\eta)\biggl\{\lambda\int_0^t\psi_0(\eta_1)\exp\Bigl(-\frac{t-\eta}{\tau_q}\Bigr)d\eta-\psi_1(\eta)\biggr\}d\eta,

\end{multline*}

$$

\bar{N}(t,\eta)=-\biggl(\frac{J_1[b(t-\eta)]}{t-\eta}+\lambda

b\int_{\eta}^tN(t,\eta_1)\frac{J_1[b(\eta_1-\eta)]}{\eta_1-\eta}d\eta_1\biggr).$$



Интегрируя равенство (\ref{el:45}) от $0$ до $t$, снова приходим к интегральному уравнению Вольтерра второго рода, ядро и правая 

часть которого непрерывны в замкнутых областях их определения. Следовательно, однозначно определяется функция $\mu(t)$ в классе 

непрерывных функций. Затем по формуле (\ref{el:39}) находим функцию $\nu^-(t)$. Решение исходной задачи в области $\Omega_2$ 

определяем по формуле (\ref{el:5}), а в области $\Omega_1$ находим по формуле (\ref{el:37}).



Пусть теперь $a\neq1$, $c\neq1$, $b=0$, $m=0$, $J\equiv(0,l_1)$. Уравнение (\ref{el:1}) будем рассматривать в области 

$\Omega=\Omega_1\bigcup\Omega_2\bigcup J$, где $\Omega_1=\{(x,t):0<x<l_1,$ $0<t<t_0\}$, $\Omega_2=\{(x,t):l_1<x<l_2,0<t\leq

t_0\}$, $l_1$, $l_2$, $t_0$ "--- константы.



\newpage



\begin{newthm}{Задача~E} Найти функцию $u(x,t),$ которая{\rm:}

\begin{enumerate}

\item[\rm 1)] является регулярным решением уравнения \eqref{el:1} в области $\Omega$ при $x\neq l_1;$

\item[\rm 2)] удовлетворяет граничным условиям

\end{enumerate}

\begin{equation} \label{el:46}

\begin{array}{c}

\displaystyle \alpha_1\frac{\partial u}{\partial x}(0,t)+\beta_1u(0,t)=\varphi_1(t),\; u(x,0)=\psi_0(x),

\\[2mm]

0\leq x\leq l_1,\; 0\leq t\leq t_0;

\end{array}

\end{equation}

\begin{equation} \label{el:47} 

\begin{array}{c}

\displaystyle

\alpha_2\frac {\partial u}{\partial x}(l_2, t)+\beta_2u(l_2,t)=\varphi_2(t),\; u(x,0)=f_0(x),\; u_t(x,0)=f_1(x),

\\[2mm]

l_1\leq x\leq l_2, \quad 0\leq t\leq t_0

\end{array}

\end{equation}

и условиям сопряжения

\begin{equation} \label{el:48}

u(-l_1,t)=u(+l_1,t), \quad \alpha_3\frac{\partial u}{\partial x}(-l_1,t)=\beta_3\frac{\partial u}{\partial x}(+l_1,t), 

\quad 0<t\leq t_0,

\end{equation}

где $\alpha_i,$ $\beta_i$ $(i=1, 2, 3),$ $l_1,$ $l_2,$ $t_0$ "--- постоянные$,$ причём $\alpha_3>0,$ $\beta_3>0,$ $\psi_0(x),$ $f_0(x)\in C^2[l_1, l_2],$ $f_1(x)\in C^1[l_1, l_2],$ $\varphi_j(t)\in C[0, t_0],$ $j=1, 2.$

\end{newthm}



Основная интегральная формула представлений решений краевых задач

для уравнения (\ref{el:1}) в области $\Omega_1$ имеет следующий вид \cite{eleevbalk:bib:Pol}:

\begin{multline} \label{el:49}

u(x,t)=\int_0^t\Bigl[\frac{\partial

u(\xi,\eta)}{\partial\xi}\Gamma(x,t;\xi,\eta)-u(\xi,\eta)\frac{\partial\Gamma(x,t;\xi,\eta)}{\partial\xi}\Bigr]_{\xi=l_2}d\eta-\\

-\int_0^t\Bigl[\frac{\partial

u(\xi,\eta)}{\partial\xi}\Gamma(x,t;\xi,\eta)-u(\xi,\eta)\frac{\partial\Gamma(x,t;\xi,\eta)}{\partial\xi}\Bigr]_{\xi=l_1}d\eta-\\

-\int_{l_1}^{l_2}\Bigl[f_1(\xi)\Gamma(x,t;\xi,0)-\frac{\partial\Gamma(x,t;\xi,\eta)}{\partial\eta}\Bigr]_{\eta=0}f_0(\xi)]d\xi,

\end{multline}

где $\Gamma(x,t;\xi,\eta)$ "--- функция Грина первой, второй или третьей краевой задачи для уравнения (\ref{el:1}) при $x>l_1$ и 

она определяется так:

$$\Gamma(x,t;\xi,\eta)=\rho\omega(x,t;\xi,\eta)+g(x,t;\xi,\eta),$$

причём

$$

{\omega(x,t;\xi,\eta)} =

\begin{cases}

h_1, & |x-\xi|\leq a(t-\eta), \\

0,   & |x-\xi|>a(t-\eta)

\end{cases}

$$

"--- функция единичного импульса,

$$

h_1= \lim_{\varepsilon\to

0}\frac{1}{2a}\int_{\xi-\varepsilon}^{\xi+\varepsilon}hdz=\frac{1}{2a}\lim_{\varepsilon\to

0}\int_{\xi-\varepsilon}^{\xi+\varepsilon}\frac{I}{2\varepsilon\rho}dz=\frac{1}{2a\rho},

$$

$\rho$ "--- линейная плотность, $g(x,t;\xi,\eta)$ "--- как функция $x$ и $t$ удовлетворяет уравнению (\ref{el:1}) при $x>l_1$ и такая, 

что $\Gamma(x,t;\xi,\eta)$ на концах рассматриваемого интервала $(l_1,l_2)$ удовлетворяет однородным граничным условиям

соответствующей краевой задачи, а $\Gamma(x,t;\xi,\eta)$ и

$\partial\Gamma(x,t;\xi,\eta)/\partial t$ равны нулю в интервале $(l_1,l_2)$ при $t<\eta$. Функция $\Gamma(x,t;\xi,\eta)$ по переменным $x$ и $t$ удовлетворяет

уравнению

$$

\Box\Gamma=\frac{\partial^2\Gamma}{\partial

x^2}-\frac{1}{a^2}\frac{\partial^2\Gamma}{\partial

t^2}=-\frac{1}{a^2}\delta(x,\xi)\delta(t,\eta),

$$

где $\delta(x,\xi)$ и $\delta(t,\eta)$ "--- $\delta$"=функции

переменных $x$ и $t$ соответственно, $\Box$ "--- оператор Даламбера. $\Gamma(x,t;\xi,\eta)$ как

функция $\xi$ и $\eta$ обладает свойством взаимности:

$\Gamma(x,t;\xi,t)=\Gamma(\xi,-\eta,x,-t)$.



В случае первой краевой задачи, т.~е. когда $\alpha_2=0$, $\beta_2=1$,

$u(+l_1,t)=\mu(t)$, представление (\ref{el:49}) принимает вид

\begin{multline*}

u(x,t)=a^2\biggl(\int_0^t[\varphi_1(\eta)\Gamma(x,t;l_1,\eta)]d\eta-\int_0^t\mu(\eta)\Gamma_{\xi}(x,t;l_2,\eta)d\eta-\\

-\int_{l_1}^{l_2}[\psi_1(\xi)\Gamma(x,t;\xi,0)-\psi_0(\xi)\Gamma_{\eta}(x,t;\xi,0)]d\xi\biggr).

\end{multline*}



В случае второй краевой задачи, т.~е.

$\alpha_2=1$, $\beta_2=0$, $u_x(-l_1,t)=(\beta_3/\alpha_3)$, $u_x(+l_1,t)=\nu^-(t)$,

представление (\ref{el:49}) имеет вид

\begin{multline*}

u(x,t)=-a^2\biggl(\int_0^t[\varphi_1(\eta)\Gamma(x,t;l_1,\eta)+\nu^-(\eta)\Gamma(x,t;l_2,\eta)]d\eta-\\

-\int_{l_1}^{l_2}[\psi_1(\xi)\Gamma(x,t;\xi,0)-\psi_0(\xi)\Gamma_{\eta}(x,t;\xi,0)]d\xi\biggr).

\end{multline*}



В случае третьей краевой задачи, т.~е. $\lambda_1u_x(\!-l_1,t)\!+\!\lambda_2u(\!-l_1,t)=\varphi_3(t)$,

представление (\ref{el:49}) принимает вид

\begin{multline*}

u(x,t)=-a^2\biggl(\int_0^t[\varphi_1(\eta)\Gamma(x,t,l_1,\eta)-\varphi_3(\eta)\Gamma(x,t;l_2,\eta)]d\eta-\\

-\int_{l_1}^{l_2}[\varphi_1(\eta)\Gamma(x,t;\xi,0)-\psi_0(\eta)\Gamma_{\eta}(x,t;\xi,0)]d\xi\biggr).

\end{multline*}



Функция Грина первой краевой задачи имеет вид \cite{eleevbalk:bib:Pol}.

$$

\Gamma(x,t;\xi,\eta)=\sum\limits_{n=1}^\infty\frac{2}{n\pi

a}\sin\frac{n\pi x}{l_1}\sin\frac{n\pi \xi}{l_1}\sin\frac{n\pi

a(t-\eta)}{l_1}.

$$



Построение функций Грина для второй, третьей и смешанных краевых задач для гиперболических уравнений в явном виде, хотя бы для 

широкого класса областей, связано со значительными трудностями. Здесь функция Грина, в отличие от эллиптических и параболических

уравнений, не приводит к сравнительно простому способу получения решения первой, второй и третьей краевых задач в явном виде хотя

бы для сколько-нибудь широкого класса областей.





\smallskip



Справедлива



\begin{theorem}[5]Пусть $\alpha_1=0,$ $\beta_1=1,$ $\alpha_2=0,$ $\beta_2=1,$ $\varphi_1(t)\in C[0, t_0];$ 

функция $\psi_0(x)$ после нечётного продолжения на отрезок $[-l_2, l_1]$ и~периодического продолжения с~периодом $2(l_2-l_1)$ принадлежит классу $C^3[l_1, l_2],$ а функция $\psi_1(x)$ при тех же условиях продолжения принадлежит классу $C^2[l_1, l_2],$ причём 

$\psi_0(l_1)=\psi_0(l_2)=0,$ $\psi''_0(l_1)=\varphi''_0(l_2)=0,$ $\psi_1(l_1)=\psi_1(l_2)=0.$ Тогда

задача $E$ имеет$,$ и притом единственное$,$ решение\rm .

\end{theorem}



\smallskip



\begin{proof}

Пусть существует решение задачи $E$, тогда, дифференцируя равенство (\ref{el:49}) по $x$, а затем переходя к пределу 

при $x\to{-l_1}$ и выполняя несложные преобразования, будем иметь

\begin{equation} \label{el:50}

\nu^-(t)=\int_0^tM(t,\eta)\mu(\eta)d\eta+m(t),

\end{equation}

где $M(t,\eta)=-a^2\Gamma_{\xi x}(-l_1,t;l_2,\eta)$, 

\begin{multline*}

m(t)=a^2\int_0^t\varphi_1(\eta)\Gamma_{\xi x}(-l_1,t;-l_2,\eta)d\eta-\\

-a^2\int_{-l_1}^{l_2}[\psi_1(\xi)\Gamma_x(-l_1,t;\xi,0)-\psi_0(\xi)\Gamma_{\eta x}(-l_1,t;\xi,0)]d\xi.

\end{multline*}



Учитывая второе условие сопряжения из (\ref{el:48}), получим 

\begin{equation} \label{el:51}

D_{0t}^{-1/2}\mu'(t)+\sqrt{\pi}\int_0^tk_0(t,\eta)\mu'(\eta)d\eta=\bar{\lambda}\int_0^tM(t,\eta)\mu(\eta)d\eta+\bar{m}(t),

\end{equation}

где

$$

\bar{\lambda}=-(a_3\sqrt{\pi})/\beta_3, \quad \bar{m}(t)=\sqrt{\pi}\psi_0(t)-(\alpha_3\sqrt{\pi})/\beta_3m(t).

$$

К обеим частям равенства (\ref{el:51}) применим оператор $D_{0t}^{+1/2}$, обратный оператору $D_{0t}^{-1/2}$. В результате будем иметь

\begin{multline*}

\mu'(t)+\sqrt{\pi}\frac{d}{dt}D_{0t}^{-1/2}\int_0^tk_0(t,\eta)\mu'(\eta)d\eta=

\\

=\bar{\lambda}\frac{d}{dt}D_{0t}^{-1/2}\int_0^tM(t,\eta)\mu(\eta)d\eta+\frac{d}{dt}D_{0t}^{-1/2}\bar{m}(t)

\end{multline*}

или

\begin{equation} \label{el:52}

\mu'(t)+\int_0^t\bar{k}_0(t,\eta)\mu'(\eta)d\eta=\bar{\lambda}\frac{d}{dt}\int_0^t\bar{M}(t,\eta)\mu(\eta)d\eta+r(t),

\end{equation}

где

$$\bar{k}_0(t,\eta)=-\frac{1}{\pi}\int_0^1z^{3/2}(1-z)^{-1/2}\mathop{\sum\rule{0mm}{11pt}{'}}\limits_{n=-\infty}^{n=\infty}\frac{n^2}{(t-\eta)^2}

\exp\Bigl(-\frac{n^2}{(t-\eta)z}\Bigr)dz,$$

$$\bar{M}(t,\eta)=\frac{1}{\sqrt{\pi}}\int_{\eta}^{t}\frac{M(\eta_1,\eta)}{(t-\eta_1)^{1/2}}d\eta_1, \quad

r(t)=\frac{d}{dt}D_{0t}^{-1/2}\bar{m}(t).$$



Легко заметить, что $\bar{k}_0(t,\eta)\in C^1[-l_1, 0]\times[-l_1, 0]\bigcap C^2]-l_1, 0[{}\times{}]-l_1, 0[$, $r(t)\in

C[0,t_0]$, $\bar{M}(t,\eta)\in C[-l_1, 0]\times[-l_1, 0]\bigcap C^1]-l_1, 0[{}\times{}]-l_1, 0[$. В уравнении (\ref{el:52}) правую часть 

пока будем считать известной. Если обозначим через $\bar{R}(t,\eta)$ резольвенту ядра $\bar{k}_0(t,\eta)$, то уравнение 

(\ref{el:51}) можно переписать в виде

\begin{multline*}

\mu(t)=r(t)+\bar{\lambda}\frac{d}{dt}\int_0^t\bar{M}(t,\eta)\mu(\eta)d\eta+\\

+\int_0^t\bar{R}(t,\eta)\biggl\{\bar{\lambda}\frac{d}{d\eta}\int_0^{\eta}\bar{M}(\eta,\eta_1)\mu(\eta_1)d\eta_1+r(\eta)\biggr\}d\eta

\end{multline*}

или

\begin{equation} \label{el:53}

\mu'(t)=\int_0^tL_1(t,\eta)\mu(\eta)d\eta+\bar{r}(t),

\end{equation}

где

$$

L_1(t,\eta)=\bar{\lambda}\bar{R}(t,t)\bar{M}(t,\eta)+\bar{M}_t(t,\eta)+\int_{\eta}^{t}\bar{R}_{\eta_1}(t,\eta_1)\bar{M}(\eta_1,\eta)d\eta_1,

$$

$$\bar{r}(t)=r(t)+\int_0^t\bar{R}(t,\eta)r(\eta)d\eta.$$



Интегрируя обе части равенства (\ref{el:53}) от $0$ до $t$, получим интегральное

уравнение Вольтерра второго рода относительно $\mu(t)$:

\begin{equation} \label{el:54}

\mu(\eta)=\int_0^t\bar{L}_1(t,\eta)\mu(\eta)d\eta+\bar{\bar{r}}(t),

\end{equation}

где

$$\bar{L}_1(t,\eta)=\int_{\eta}^tL_1(\eta_1,\eta)d\eta_1, \quad \bar{\bar{r}}(t)=\int_0^t\bar{r}(\eta)d\eta.$$



На основании свойств заданных функций

$\varphi_1(t)$, $\varphi_2(t)$, $\psi_0(x)$, $\psi_1(x)$, $f(x)$ заключаем, что

$\bar{L}_1(t,\eta)\in C[-l_1,0]\times[-l_1,0]$, $\bar{\bar{r}}(t)\in

C[0,t_0]$. Следовательно, уравнение (\ref{el:54}) безусловно и однозначно

разрешимо, и его решение $\mu(t)\in C[0,t_0]$.



После определения функции $\mu(t)$ решение задачи $E$ в области

$\Omega_2$ находим по формуле (\ref{el:49}), а в области $\Omega_1$ оно

задаётся так:

\begin{multline*}

u(x,t)=\int_0^{+l_1}f(\xi)G(\xi,0;x,t)d\xi+\int_0^t\mu(\eta)G_\xi(0,\eta;x,t)d\eta-\\

-\int_0^t\varphi_1(\eta)G_\xi(+l_1,\eta;x,t)d\eta,

\end{multline*}

где $G(\xi,\eta;x,t)$ "--- функция Грина первой краевой задачи уравнения (\ref{el:1}) при $x<l_1$ и определяется формулой (\ref{el:7}).

\end{proof}



\begin{thebibliography}{99}



\RBibitem{eleevbalk:bib:Gel}

\by Гельфанд~И.\,И.

\paper Некоторые вопросы анализа и дифференциальных уравнений

\jour УМН

\yr 1959

\vol 14

\issue 3(87)

\pages 3--19

%\zmath{http://www.zentralblatt-math.org/zmath/search/?an=Zbl 0131.31802|Zbl 0091.08805}

\transl англ. пер.:~

\by Gel'fand~I.\,M.

\paper Some questions of analysis and differential equations

\jour Am. Math. Soc., Transl., II.~Ser. 

\vol 26

\pages 201--219 

\yr 1963



\RBibitem{eleevbalk:bib:Leib}

\by Лейбензон~Л.\,Л.

\book Движение природных жидкостей и газов в пористой среде

\publaddr М.,~Л.

\publ ОГИЗ, Гостехиздат

\yr 1947

\totalpages 244

\misctransl

\by Leybenzon~L.\,L.

\book  Motion of natural liquids and gases in porous medium

\publaddr Moscow, Leningrad

\publ OGIZ, Gostekhizdat

\yr 1947

\totalpages 244



%3

\RBibitem{eleevbalk:bib:TixSam}

\book Уравнения математической физики

\by Тихонов~А.\,Н., Самарский~А.\,А.

\publaddr М.

\publ Наука

\yr 1977

\totalpages 735

\transl англ. пер.:~

\by Tikhonov~A.\,N., Samarskii~A.\,A.

\book Equations of Mathematical Physics 

\publ Pergamon Press

\publaddr Oxford

\yr 1963

\totalpages 784





%4

\RBibitem{eleevbalk:bib:Aziz}

\book Математическое моделирование пластовых систем

\by  Азиз~Х., Сеттари~Э.

\publaddr М.

\publ Недра

\yr 1982

\totalpages 407

\misctransl

\by Aziz~H., Settari~E.

\book Mathematical modeling of reservoir systems

\publaddr Moscow

\publ Nedra

\yr 1982

\totalpages 407



%5

\RBibitem{eleevbalk:bib:Ak}

\book Нестационарные движения вязкоупругих жидкостей

\by  Акилов~Ж.\,А.

\publaddr Ташкент

\publ Фан

\yr 1982

\totalpages 104

\misctransl

\by Akilov~J.\,A.

\book Non-stationary motions of viscoelastic fluids

\publ Fan

\publaddr Tashkent

\yr 1982

\totalpages 104





%6

\RBibitem{eleevbalk:bib:Kirz}

\paper Разрешимость задачи сопряжения гиперболического и параболического уравнений с~интегро"=дифференциальными условиями на границе раздела областей

\by Корзюк~В.\,И., Лемешевский~С.\,В., Матус~П.\,П.

\jour Тр. ин-та мат. НАН Беларуси

\yr 2000

\vol 6

\pages 100--108

\misctransl

\by Korzyuk~V.\,I., Lemeshevskiy~S.\,V., Matus~P.\,P.

\paper Solvability of a conjugation problem for hyperbolic and parabolic equations with integro-differential conditions on the separation boundary

\jour Tr. in-ta mat. NAN Belarusi

\yr 2000

\vol 6

\pages 100--108



%7

\RBibitem{eleevbalk:bib:Shash}

\book Системно-структурный анализ процесса теплообмена и его применение

\by  Шашков~А.\,Г.

\publaddr М.

\publ Энергоатомиздат

\yr 1983

\totalpages 279

\misctransl

\by Shashkov~A.\,G.

\book Systematic Structural Analysis of Heat Transfer Processes and its Application

\publ Energoizdat

\publaddr Moscow 

\yr 1983

\totalpages 279



%8

\RBibitem{eleevbalk:bib:Chudn}

\book Теплофизика почв

\by  Чудновский~А.\,Ф.

\publaddr М.

\publ Наука

\yr 1976

\totalpages  352

\misctransl

\by Chudnovskiy~A.\,F.

\book Thermophysics of Soils 

\publ Moscow 

\yr 1976

\publ Nauka

\totalpages  352



%9

\RBibitem{eleevbalk:bib:Str}

\paper Задача о сопряжении двух уравнений

\by Стручина~Г.\,М.

\jour Инженерно-физический журнал

\yr 1961

\vol 4

\issue 11

\pages 99--104

\misctransl

\by Struchina~G.\,M.

\paper A problem of conjugation of two equations

\jour Inzherno-Fizicheskii Zhurnal 

\vol 4

\issue 11

\pages 99–104

\yr 1961



%10

\RBibitem{eleevbalk:bib:Zol}

\by Золина~Л.\,А.

\paper О краевой задаче для модельного уравнения гиперболо-параболического типа

\jour Ж.~вычисл. матем. и матем. физ.

\yr 1966

\vol 6

\issue 6

\pages 991--1001

\transl англ. пер.:~

\paper On a boundary value problem for a model equation of hyperbolo-parabolic type

\by Zolina~L.\,A.

\jour U.S.S.R. Comput. Math. Math. Phys.

\yr 1966

\vol 6

\issue 6

\pages 63--78



%11

\RBibitem{eleevbalk:bib:Abdr}

\paper Об одной задаче сопряжения уравнений параболического типов

\by Абдрахманов~М.\,А.

\jour Извест. АН КазССР. Сер. физ.-мат. науки

\yr 1967

\vol 5

\pages 87--93

\misctransl 

\by Abrakhmanov~M.\,A.

\paper On one  conjugation problem for parabolic type equations 

\jour Izvest. AN KazSSR. Ser. fiz.-mat. nauki

\yr 1967

\vol 5

\pages 87--93



%12

\RBibitem{eleevbalk:bib:Ostr}

\paper Задача на сопряжение уравнений параболического и гиперболического типов, когда в граничные условия входят производные 

по времени

\by Островский~Е.\,А.

\jour Дифференц. уравнения

\yr 1967

\vol 3

\issue 6

\pages 965--979

\misctransl

\by Ostrovskiy~E.\,A.

\jour Differents. uravneniya

\yr 1967

\vol 3

\issue 6

\pages 965--979



%13

\RBibitem{eleevbalk:bib:Kor}

\paper Задача о сопряжении уравнений гиперболического и параболического типов

\by Корзюк~В.\,И.

\jour Дифференц. уравнения

\yr 1968

\issue 4

\pages 1855--1866

\misctransl

\by Korzyuk~V.\,I.

\paper The problem of conjugate equations of hyperbolic and parabolic type

\jour Differents. uravneniya

\yr 1968

\issue 4

\pages 1855--1866





\RBibitem{eleevbalk:bib:Jur}

\book Краевые задачи для уравнений смешанного и смешанно-составного типа.

\by Джураев~Т.\,Д.

\publaddr Ташкент

\publ Фан

\yr 1979

\totalpages 238

\misctransl

\by Dzhuraev~T.\,D.

\book Boundary value problems for equations of mixed and mixed-composite type

\publaddr Tashkent 

\yr 1979

\totalpages 238

\publ Fan





\RBibitem{eleevbalk:bib:JurSop}

\book Краевые задачи для уравнений параболо-гиперболического типа

\by Джураев~Т.\,Д., Сопуев~А.\,C., Мамажанов~М.

\publaddr Ташкент

\publ Фан

\yr 1986

\totalpages 220

\misctransl 

\by Dzhuraev~T.\,D., Sopuev~A.\,S., Mamazhanov~M.

\book Boundary value problem for equations of hyperbolic-parabolic type

\publaddr Tashkent

\publ Fan

\yr 1986

\totalpages 220





\RBibitem{eleevbalk:bib:Baz}

\thesis К теории локальных и нелокальных краевых задач для уравнений смешанного и смешанно"=составного типов

\thesisinfo Дис. \dots\ д-ра физ.-мат. наук

\by Базаров~Д.\,К.

\publaddr Ашхабад

\yr 1991

\totalpages 272

\misctransl

\by Bazarov~D.\,K.

\thesis For the theory of local and nonlocal boundary value problems for equations of mixed and mixed-composite type

\publ Ashgabat

\thesisinfo Dr. Sci. Thesis

\yr 1991

\totalpages 272





\RBibitem{eleevbalk:bib:El}

\thesis Краевые задачи для уравнений смешанного гиперболо"=параболического типа

\thesisinfo Дис. \dots\ д-ра физ.-мат. наук

\by Елеев~В.\,А.

\publaddr Киев

\yr 1995

\totalpages 266

\misctransl 

\by Eleev~V.\,A.

\thesis Boundary value problems for equations of mixed hyperbolic-parabolic type

\thesisinfo Dr. Sci. Thesis

\publaddr Kiev

\yr 1995

\totalpages 266



\RBibitem{eleevbalk:bib:Rep}

\thesis Краевые задачи для уравнений гиперболического и смешанного типов и дробное интегро"=дифференцирование

\thesisinfo Дис. \dots\ д-ра физ.-мат. наук

\by Репин~О.\,А.

\publaddr Минск

\yr 1998

\totalpages 220

\misctransl

\by Repin~O.\,A.

\thesis Boundary value problems for equations of hyperbolic and mixed types and fractional integro-differentiation

\thesisinfo Dr. Sci. Thesis

\publaddr Minsk

\yr 1998

\totalpages 220







\RBibitem{eleevbalk:bib:Sab}

\thesis Некоторые вопросы качественной спектральной теории уравнений смешанного типа

\thesisinfo Дис. \dots\ д-ра физ.-мат. наук

\by Сабитов~К.\,Б.

\publaddr M.

\yr 1991

\totalpages 313

\misctransl

\by Sabitov~K.\,B.

\thesis Some problems in qualitative and spectral theory of mixed-type equations 

\thesisinfo Dr. Sci. Thesis

\publaddr Moscow 

\yr 1991

\totalpages 313





\RBibitem{eleevbalk:bib:Nah}

\by Нахушева~В.\,А.

\paper Об одной математической модели теплообмена в~смешанной среде с~идеальным контактом

\jour  Вестн. Сам. гос. техн. ун-та. Сер. Физ.-мат. науки

\yr 2006

\issue 42

\pages 11--34

\misctransl

\by Nakhusheva~V.\,A.

\paper On one mathematical model of heat transfer in a mixed media with ideal contact

\jour Vestn. Samar. Gos. Tekhn. Univ. Ser. Fiz.-Mat. Nauki

\yr 2006

\issue 42

\pages 11--34





\RBibitem{eleevbalk:bib:Eleev}

\paper О краевых задачах для смешанного уравнения гиперболо-параболического типа

\by Елеев~В.\,А.

\jour Дифференц. уравнения

\yr 1988

\vol 24

\issue 4

\pages 627--635

\transl англ. пер.:~

\by Eleev~V.\,A.

\paper Boundary-value problems for a mixed hyperbolic-parabolic equation

\jour Differ. Equations 

\vol 24

\issue 4

\pages 437--443 

\yr 1988







\RBibitem{eleevbalk:bib:Mis}

\book Математика для ВТУЗов

\bookinfo Специальные курсы

\by Мышкис~А.\,Д.

\publaddr M.

\publ Наука

\yr 1971

\totalpages 632

\misctransl

\by Myshkis~A.\,D.

\book Mathematics for technical universities

\bookinfo Special courses

\publaddr Moscow

\publ  Nauka

\yr 1971

\totalpages 632





\RBibitem{eleevbalk:bib:Salah}

\book Краевые задачи для уравнений смешанного типа со спектральным параметром

\by Салахитдинов~М.\,С., Уринов~А.\,К.

\publaddr Ташкент

\yr 1997

\publ Фан

\totalpages 165 

\misctransl

\by Salakhitdinov~M.\,S., Urinov~A.\,K.

\book Boundary value problems for the mixed type equations with spectral parameter

\publ Fan

\publaddr Tashkent

\yr 1997

\totalpages 166





\RBibitem{eleevbalk:bib:Pol}

\book Уравнения математической физики

\by Положий~Г.\,П.

\publaddr M.

\publ Высш. шк.

\yr 1964

\totalpages 559

\misctransl

\by Polozhiy~G.\,P.

\book Equations of mathematical physics

\publaddr Moscow

\publ Vyssh. shk.

\yr 1964

\totalpages 559



\end{thebibliography}



\noindent

\hskip6.5cm \thedate \par\noindent

\hskip6.5cm\therevdate



\vspace{-2mm}



\makeengtitle



\newpage

%

%\end{document}